\begingroup
\renewcommand{\arraystretch}{1.45}
\fancyhead[L]{\small BESTIMMUNG DER ANZIEHUNG}
\Needspace{42\baselineskip}
\AttractionArticle{10}

Unsere Formeln können in gewissen Fällen unbestimmt werden, die gesondert betrachtet werden müssen. Untersuchen wir zunächst den Fall, in dem die negativen Wurzeln $-G'$, $-G''$ der kubischen Gleichung gleich werden, sodass die Koeffizienten $\gamma'$, $\gamma''$ durch die Formeln 18, 19 unendliche Werte zu erhalten scheinen, die aber tatsächlich unbestimmt sind.

Setzt man in Formel 14 $g=G$, so ist ersichtlich, dass, damit die beiden Werte von $x$, das heißt $-G'$ und $-G''$, gleich werden, notwendig gelten muss
\[
\begin{aligned}
&AB=0,\quad aa-bb-\frac{aaAA}{aa+G}+\frac{bbBB}{bb+G}=0
\end{aligned}
\]
\par Da $aa-bb$ seiner Natur nach entweder eine positive Größe oder $=0$ ist, erkennt man hieraus leicht, dass gelten muss
\[
\begin{aligned}
&B=0
\end{aligned}
\]
\[
\begin{aligned}
&aa-bb=\frac{aaAA}{aa+G},\quad\text{oder}\quad aa+G=\frac{aaAA}{aa-bb}
\end{aligned}
\]
Setzt man diese Werte in Gleichung 14 ein, so ergibt sich
\[
\begin{aligned}
&G'=G''=bb
\end{aligned}
\]
Setzt man ferner den Wert $x=-bb$ in die kubische Gleichung 13 ein, so ergibt sich
\[
\begin{aligned}
&(aa-bb)(CC+bb)=bbAA
\end{aligned}
\]
Sooft diese Bedingungsgleichung zugleich mit der Gleichung $B=0$ gilt, tritt der hier behandelte Fall ein. \Needspace{31\baselineskip}
Und da
\[
\begin{aligned}
&G=\frac{aaAA}{aa-bb}-aa=\frac{aaCC}{bb}
\end{aligned}
\]
\noindent\SourceLeaf{356}
gilt, liefert Formel 17
\[
\begin{aligned}
&\gamma=\surd{\frac{aabbAA}{(aa-bb)(aaCC+b^4)}}
=\surd{\frac{aaCC+aabb}{aaCC+b^4}}
\end{aligned}
\]
und dann liefern die Formeln 3, 4
\[
\begin{aligned}
&\alpha=\frac{\gamma(aa-bb)}{aA}=\frac{\gamma bbA}{a(CC+bb)}
=\surd{\frac{bb(aa-bb)}{aaCC+b^4}}
\\[6pt]
&\qquad{}=\surd{\frac{b^4AA}{(CC+bb)(aaCC+b^4)}}
\end{aligned}
\]
\[
\begin{aligned}
&\beta=0
\end{aligned}
\]
\Needspace{29\baselineskip}
\par Die Werte der Koeffizienten $\gamma'$, $\gamma''$ bleiben durch die Formeln 18, 19 in diesem Fall unbestimmt, und ebenso die Werte der übrigen Koeffizienten $\alpha'$, $\beta'$, $\alpha''$, $\beta''$. Dennoch lassen sich durch einen dieser Koeffizienten alle fünf übrigen ausdrücken; zum Beispiel ergibt sich aus Formel 6
\[
\begin{aligned}
&\alpha'=\frac{\gamma'aA}{aa-bb}
\end{aligned}
\]
und dann
\[
\begin{aligned}
&\beta'=\surd{(1-\alpha'\alpha'+\gamma'\gamma')},\\[6pt]
&
\gamma''=\surd{(\gamma\gamma-1-\gamma'\gamma')},\\[6pt]
&
\alpha''=\frac{\gamma''aA}{aa-bb},\\[6pt]
&
\beta''{}=\surd{(1-\alpha''\alpha''+\gamma''\gamma'')}
\end{aligned}
\]
\Needspace{19\baselineskip}
Eleganter wird dies aber auf folgende Weise ausgeführt. Aus
\[
\begin{aligned}
&\gamma\gamma=1+\alpha\alpha,\quad
\alpha\alpha'=\gamma\gamma',\quad
1=\alpha'\alpha'+\beta'\beta'-\gamma'\gamma'
\end{aligned}
\]
folgt
\[
\begin{aligned}
&\beta'\beta'+\frac{\gamma'\gamma'}{\alpha\alpha}
=1-\alpha'\alpha'+\frac{\gamma\gamma\gamma'\gamma'}{\alpha\alpha}=1
\end{aligned}
\]
\Needspace{32\baselineskip}
\par Daher können wir setzen
\[
\begin{aligned}
&\beta'=\cos f,\quad\gamma'=\alpha\sin f,\quad\alpha'=\gamma\sin f
\end{aligned}
\]
Dann finden wir aus den Formeln
\[
\begin{aligned}
&\varepsilon\alpha''=\beta\gamma'-\gamma\beta',\quad
\varepsilon\beta''=\gamma\alpha'-\alpha\gamma',\quad
\varepsilon\gamma''=\beta\alpha'-\alpha\beta',\quad
\varepsilon\varepsilon=1
\end{aligned}
\]
\[
\begin{aligned}
&\alpha''=-\varepsilon\gamma\cos f,\quad
\beta''=\varepsilon\sin f,\quad
\gamma''=-\varepsilon\alpha\cos f
\end{aligned}
\]
Der Wert des Winkels $f$ ist hier beliebig, und ebenso kann nach Belieben entweder $\varepsilon=+1$ oder $\varepsilon=-1$ gesetzt werden.

\Needspace{43\baselineskip}
\AttractionArticle{11}

Sind $G'$, $G''$ ungleich, so können die Werte der Koeffizienten $\gamma$, $\gamma'$, $\gamma''$ durch die Formeln 17, 18, 19 nicht unbestimmt sein. Sooft aber eine der Größen
\noindent\SourceLeaf{357}
$aa-G'$, $bb-G'$, $aa-G''$, $bb-G''$ verschwindet, scheint der Wert des Koeffizienten $\alpha'$, $\beta'$, $\alpha''$, $\gamma''$ durch die Formel 6, 7, 9 beziehungsweise 10 auf den ersten Blick unbestimmt zu bleiben; ein wenig Aufmerksamkeit wird jedoch zeigen, dass es sich anders verhält.

Nehmen wir zum Beispiel an, dass $aa-G'=0$ gilt; dann wird nach Gleichung 18 $\gamma'=0$ und ebenso nach Gleichung 7 $\beta'=0$ (sofern nicht zugleich $aa=bb$ gilt), woraus notwendig $\alpha'=\pm1$ folgen muss. Gilt aber zugleich $aa=bb$, so liefert die Formel, die in Art. 5 der sechsten vorangeht, $\alpha'A+\beta'B=0$; diese Gleichung ergibt zusammen mit $\alpha'\alpha'+\beta'\beta'=1$
\[
\begin{aligned}
&\alpha'=\frac{B}{\surd{(AA+BB)}},\quad
\beta'=\frac{-A}{\surd{(AA+BB)}}
\end{aligned}
\]
\par Diese Ausdrücke können offenbar nicht unbestimmt sein, es sei denn, es gilt zugleich $A=0$, $B=0$; dann werden wir aber auf den bereits im vorangehenden Artikel betrachteten Fall zurückkommen.

\clearpage
\endgroup
